\documentclass[12pt]{article}
\usepackage[utf8]{inputenc}
\usepackage[T1]{fontenc}
\usepackage{hyperref}
\title{Choogen Language Resources for K-12 Curriculum Development}
\author{Compiled by Claude Sonnet 5 (claude-sonnet-5) for Andr\'as Kornai}
\date{}
\begin{document}
\maketitle

\section{Introduction}
Choogen (ISO 639-3 tgd, Glottocode ciwo1236) is an Afro-Asiatic (West Chadic) language of Bauchi State, Nigeria. In the sources found, it appears mainly under the name Ciwogai (also Tsagu, Sago, Tsaganci); ``Choogen Tsagu'' is listed by Joshua Project as a dialect label. Speaker estimates differ: about 2,000 (Wikipedia, 1995) versus 4,800 (Joshua Project). No dedicated educational material was found.

\section{Primary Educational Resources}
None found. No primer, textbook or school programme for this language turned up in search results.

\section{Linguistic and Reference Materials}
\begin{itemize}
\item Ciwogai language, Wikipedia overview: \url{https://en.wikipedia.org/wiki/Ciwogai_language}
\item Glottolog entry (classification, references): \url{https://glottolog.org/resource/languoid/id/ciwo1236.bigmap.html}
\item Numerals database entry (traditional numbers only 1--5 recorded): \url{https://lingweb.eva.mpg.de/channumerals/Ciwogai.htm}
\item ASJP source record (word-list data): \url{https://asjp.clld.org/sources/2506}
\end{itemize}

\section{Religious and Community Resources}
Joshua Project reports no Bible translation or scripture and an unknown written status: \url{https://joshuaproject.net/languages/tgd}. No community organisation was found.

\section{Digital and Audio Resources}
None found.

\section{Curriculum Development Context}
The language is essentially undocumented for educational purposes. Any curriculum would first need orthography development, community consultation and basic documentation (lexicon, texts, audio). The name Choogen should be reconciled with Ciwogai/Tsagu with the community before use.

\end{document}